% ============================================
% 第1章 绪论
% ============================================
\setcounter{page}{1}
\pagenumbering{arabic}

\section{绪论}

\subsection{研究背景与意义}

随着数字化转型的深入推进，数据驱动决策已成为各行各业的基本共识。问卷调查作为一种经典的数据收集方法，在互联网技术的赋能下焕发出新的活力。在线问卷系统以其低成本、高效率、广覆盖的优势，广泛应用于市场调研、学术研究、用户反馈、满意度评估等诸多场景。

当前，国内外已涌现出众多在线问卷平台，如国外的 SurveyMonkey、Google Forms，国内的问卷星、腾讯问卷等。这些平台功能成熟、用户基数庞大，但在实际使用中也暴露出一些共性问题：一是协作能力薄弱，多人共同编辑问卷时缺乏实时同步机制；二是智能化程度不足，问卷设计和数据分析仍高度依赖人工经验；三是定制化能力有限，难以满足特定场景下的个性化需求。

在此背景下，本文设计并实现了一个基于 Next.js 的在线问卷系统。系统不仅涵盖了问卷创建、发布、收集、统计等完整生命周期管理，还创新性地引入了 AI 辅助生成与分析和多人实时协作编辑功能，旨在为用户提供一个更加智能、高效的问卷工作平台。

\subsection{国内外研究现状}

\subsubsection{国外研究现状}

国外在线问卷市场起步较早，发展较为成熟。SurveyMonkey 成立于 1999 年，是全球最大的在线问卷平台之一，提供丰富的题型模板和数据分析工具。Google Forms 依托 Google 生态，以简洁易用著称，并与 Google Sheets 深度集成。Typeform 则以交互式问卷设计见长，注重用户体验和视觉呈现。

在学术研究领域，问卷系统的技术演进也备受关注。近年来，随着自然语言处理技术的发展，研究者开始探索利用 AI 辅助问卷设计。例如，有研究利用大语言模型自动生成调查问题，并通过迭代优化提升问题质量。在协作方面，Operational Transformation（OT）和 Conflict-free Replicated Data Types（CRDT）等算法为实时协作编辑提供了理论基础。

\subsubsection{国内研究现状}

国内在线问卷市场同样发展迅速。问卷星作为国内老牌问卷平台，拥有庞大的用户群体和丰富的功能模块。腾讯问卷依托微信生态，在社交传播方面具有天然优势。金数据、麦客等新兴平台则在表单设计和数据收集方面各有特色。

在学术研究方面，国内学者主要关注问卷系统的可用性优化、移动端适配、数据安全等方面。近年来，随着国产大模型的崛起，AI 与问卷系统的结合也成为新的研究方向。

\subsubsection{现有平台的不足}

综合分析现有平台，主要存在以下不足：

\begin{enumerate}[label=(\arabic*)]
  \item \textbf{协作能力有限}：多数平台仅支持简单的问卷分享，缺乏真正的多人实时协作编辑能力。
  \item \textbf{AI 集成浅显}：部分平台虽引入了 AI 功能，但多停留在简单的文本生成层面，缺乏深度的数据分析和洞察能力。
  \item \textbf{版本管理缺失}：问卷内容频繁修改时，缺乏有效的版本回溯机制。
  \item \textbf{定制化不足}：平台功能固定，难以根据特定需求进行灵活扩展。
\end{enumerate}

\subsection{论文主要工作}

本文的主要工作包括以下几个方面：

\begin{enumerate}[label=(\arabic*)]
  \item \textbf{系统架构设计}：基于 Next.js 全栈框架，设计前后端一体化的系统架构，采用 Server Components 优化渲染性能。
  
  \item \textbf{问卷编辑器实现}：开发可视化问卷编辑器，支持 18 种题型，提供拖拽排序、实时预览等功能。
  
  \item \textbf{版本管理机制}：设计并实现问卷版本管理系统，支持发布快照、版本回滚等功能。
  
  \item \textbf{AI 辅助功能}：集成 DeepSeek 大语言模型，实现问卷智能生成和数据分析总结功能。
  
  \item \textbf{实时协作编辑}：基于 Pusher 实现多人实时协作编辑，支持题目级乐观锁定和状态同步。
  
  \item \textbf{数据可视化}：利用 ECharts 实现多维度的数据统计和可视化展示。
\end{enumerate}

\subsection{论文组织结构}

本文共分为六章，各章内容安排如下：

\textbf{第1章 绪论}：介绍研究背景与意义，分析国内外研究现状，阐述论文主要工作和组织结构。

\textbf{第2章 相关技术介绍}：介绍系统开发所涉及的关键技术，包括 Next.js、React、Prisma、Pusher、AI SDK 等。

\textbf{第3章 系统需求分析与设计}：进行需求分析，设计系统架构、数据库模型和接口规范。

\textbf{第4章 系统详细设计与实现}：详细阐述各核心模块的设计与实现过程，包括编辑器、发布、填写、统计、AI、协作等模块。

\textbf{第5章 系统测试}：对系统进行功能测试、性能测试和兼容性测试，分析测试结果。

\textbf{第6章 总结与展望}：总结论文工作，分析系统不足，展望未来改进方向。
